\input format.tex

\usepackage{graphicx}
\graphicspath{{cores/}}

\begin{document}

\vspace*{3mm}
\parindent0pt
%% 各章节
\setlength{\arrayrulewidth}{1pt}
\fontsize{9.03pt}{14pt}\selectfont
\color{gray2}

{\bf\biaotia 检测结果总览}

\vspace*{2mm}

%\begin{spacing}{1.25}
\begin{LRaside}{肠道菌群概况}
\noindent\centering
\includegraphics[width=\linewidth]{gaikuang_red.pdf}
\jiuhao
\noindent\begin{tabular}{@{}m{0.68cm}<{\centering}@{}m{0.68cm}<{\centering}@{}m{0.68cm}<{\centering}@{}m{0.68cm}<{\centering}@{}}
\includegraphics[width=.41cm]{sun_red.pdf} & \tikz\draw[gray,fill=gray](0,0) circle[radius=4.1pt]; & \tikz\draw[gray,fill=gray](0,0) circle[radius=4.1pt]; & \tikz\draw[gray,fill=gray](0,0) circle[radius=4.1pt]; \\[-4pt]
\color{gray2}当次 & \color{gray2} & \color{gray2} & \color{gray2}
\\
\end{tabular}
\asidebreak % 分栏
1、您本次肠道菌群共检测出细菌种类175种，肠道菌群多样性
偏低，说明您的肠道菌群组成较为单一，菌群失调风险较高。
您肠道中的有益菌含量
偏低，仅高于6.77{\%}的参考人群。
有害菌含量
与参考人群相当。\par\medskip
2、您肠道内的
柔嫩梭菌属、罗斯拜瑞氏菌属、萨特氏菌属、弯曲杆菌属4项指标异常，不利于肠道健康，需引起注意，建议持续监测肠道健康。\par\medskip
3、您的肠道内
检测出4种致病菌，含量均未超标，但在您免疫力低下时仍可能会损害您的肠道健康。
\end{LRaside}

\begin{LRaside}{代谢分析及免疫}
\noindent
\includegraphics[width=\linewidth]{fenbu_green.pdf}

\jiuhao
\noindent\begin{tabular}{@{}m{0.68cm}<{\centering}@{}m{0.68cm}<{\centering}@{}m{0.68cm}<{\centering}@{}m{0.68cm}<{\centering}@{}}
\includegraphics[width=.41cm]{sun_green.pdf} & \tikz\draw[gray,fill=gray](0,0) circle[radius=4.1pt]; & \tikz\draw[gray,fill=gray](0,0) circle[radius=4.1pt]; & \tikz\draw[gray,fill=gray](0,0) circle[radius=4.1pt]; \\[-4pt]
\color{gray2}当次 & \color{gray2} & \color{gray2} & \color{gray2}
\\
\end{tabular}
\asidebreak % 分栏
1、您肠道内的
主要营养物质代谢指标均处于正常范围，有利于肠道菌群及功能的稳定，帮助维持肠道健康。\par\medskip
2、
 %只有有益
您肠道内参与维生素B6、维生素B7代谢的菌群含量优于参考人群，有益于人体健康，请保持良好的饮食及生活习惯。\par\medskip
3、
您肠道菌群的免疫指数指标处于正常范围，肠道屏障功能正常。
\end{LRaside}


\begin{LRaside}{相关疾病风险分析}
\noindent
\includegraphics[width=\linewidth]{gongneng_orange.pdf}

\jiuhao
\noindent\begin{tabular}{@{}m{0.68cm}<{\centering}@{}m{0.68cm}<{\centering}@{}m{0.68cm}<{\centering}@{}m{0.68cm}<{\centering}@{}}
\includegraphics[width=.41cm]{sun_orange.pdf} & \tikz\draw[gray,fill=gray](0,0) circle[radius=4.1pt]; & \tikz\draw[gray,fill=gray](0,0) circle[radius=4.1pt]; & \tikz\draw[gray,fill=gray](0,0) circle[radius=4.1pt]; \\[-4pt]
\color{gray2}当次 & \color{gray2} & \color{gray2} & \color{gray2}
\\
\end{tabular}
\asidebreak %分栏
综合本次肠道菌群检测结果，
提示会增加炎症性肠病的风险，需尽量避免相关疾病的高危因素，注意调理肠道菌群结构，改善肠道环境，降低疾病风险，还需注意定期检查。
\end{LRaside}
%\end{spacing}

{\bf 颜色指示}\hspace*{2mm}\includegraphics[scale=1]{greencircle.pdf} {\bf 绿色表示健康}\hspace*{2mm}\includegraphics[scale=1]{orangecircle.pdf}  {\bf 橙色表示需要关注}\hspace*{2mm}\includegraphics[scale=1]{redcircle.pdf} {\bf 红色表示有风险}


\end{document}
